\section{Choose accumulator widths from a bound}
A sample is in $[0,255]$, so a difference magnitude is at most 255. One full row
can contribute $16\times255=4080$, and a full block can contribute
$256\times255=65\,280$. A signed halfword cannot hold that final total. A byte
accumulator can wrap even sooner: sixteen contributions of 255 sum to 4080,
whose low byte is 240.

For the vector path, take unsigned absolute differences in byte lanes. Then widen
pairs before accumulating across rows. One adjacent pair contributes at most 510
per row; across sixteen rows, its halfword lane is bounded by 8160. Eight such
lanes cover the whole block. Reducing them through wider word and doubleword
lanes avoids losing the carries needed for the final sum.

\begin{table}[htbp]
\centering
\caption{Bounds for the chosen reduction. A lane's interpretation changes at each widening.}
\tagpdfsetup{table/header-rows={1}}
\begin{tabular}{lrr}
\toprule
Quantity & Lanes & Maximum per lane\\
\midrule
One row's byte differences & 16 & 255\\
One row's adjacent-pair sums & 8 & 510\\
Sixteen accumulated rows & 8 & 8160\\
Word reduction & 4 & 16\,320\\
Doubleword reduction & 2 & 32\,640\\
Final scalar total & 1 & 65\,280\\
\bottomrule
\end{tabular}
\end{table}

Using the same source twice for \texttt{hadd\_u.h} combines the adjacent byte
values into halfwords. Accumulate those halfwords with \texttt{addv.h}. After the
sixteen rows, reduce with \texttt{hadd\_u.w} and \texttt{hadd\_u.d}, extract the
two doubleword lanes and add them in a general register. These operations are
integer operations; a floating-point rounding mode is not part of the SAD metric.

For example, with row samples already in \texttt{\$w0} and \texttt{\$w1},
and a previously zeroed halfword accumulator in \texttt{\$w4}, the row contribution is:
\begin{minted}[samepage=true]{mars-idea}
    asub_u.b $w2, $w0, $w1
    hadd_u.h $w3, $w2, $w2
    addv.h $w4, $w4, $w3
\end{minted}
The first instruction keeps sixteen magnitudes in byte lanes. The second produces
eight pair sums, and the third retains contributions from earlier rows. The vector
register numbers are choices for this example, not requirements of the interface.

Both input vectors undergo the same lane mapping. Reversing the lane order
consistently does not change the final sum, but reading an extra byte or switching
to signed sample differences does. Test zero against 255: unsigned difference
magnitudes are 255, while interpreting the latter byte as signed minus one would
produce the wrong metric. All-zero, all-maximum and equal-window cases are useful
because a wrapped result can look plausible on ordinary pictures.
